\subsection{segmentation}

\begin{itemize}

\item Segmentation d'images

La segmentation appliquée à l'image en général, y compris l'image fixe, consiste à diviser l'image en groupes de pixels selon certains critères. 

\item Segmentation vidéo

Appliquée à l'image animée, la segmentation peut vouloir dire plusieurs choses différentes. En effet, l'image audiovisuelle comporte une valeur ajoutée à l'image fixe : celle de la dimension temporelle d'une succession d'images fixes créant une animation de l'image. Elle peut aussi comporter une piste audio, dans la plupart des cas, puisque le son est intégré au cinéma dès 1932. Il faut tout de même noter que toutes les images animées ne comportent pas de son (films de famille, films expérimentaux, films muets). Dans notre analyse, le son sera tout de même pris en compte dans la mise en place d'un pipeline de segmentation pour prévoir le cas fréquent d'une bande sonore. 

Une fois ceci expliqué, la segmentation peut être  appliquée à l'image audiovisuelle de manières différentes. Généralement, elle peut être définie comme une division d'une vidéo en plusieurs segments ou objets.\footcite{zhou2022} 
Plus spécifiquement, elle peut désigner la segmentation d'une séquence en différents plans, ou d'un film en différentes séquences, on parlera alors de segmentation temporelle. 
Elle peut aussi désigner la segmentation des objets "dominants" \footcite{zhou2022 }dans l'image selon des critères visuels.\footcite{carrive2014} Ce type de segmentation cherche seulement à analyser la place des objets dans l'image, et non leur signification. \footcite{zhou2022}. 
On parle de segmentation sémantique lorsque que survient une extraction des objets qui apparatiennent à des catégories sémantqiues prédéfinies. \footcite{zhou2022} AInsi la segmentation sémantique prépare l'indexation en assignant des étiquettes à chaque pixel extrait de l'image. Elle est aussi applicable à l'image fixe, puisqu'elle ne prend pas en compte la dimension temporelle de l'image anme, contrairement à la segmentation temporelle. 
Enfin, peut être appliquée à une vidéo une segmentation de la bande sonore, qu'on désignera sous le nom de segmentation audio, par exemple entre des paroles et de la musique. \footcite{carrive2014}

\item la segmentation sémantique 
Concentrons nous ici sur la segmentation sémantique. Deux modèles de segmentation sémantique se distinguent. Les méthodes orientées pécision, qui agrgègent l'information temporelle entre les images our améliorer la prédiction. (Elles utilisent le flux optique ou des modles basés sur l'attention. Ex de méthodes : TDNET) \footcite{zhou2022}
Les méthodes orientées vitesse, qui réutilisent les caractéristiquesdes images adjacentes pour ne pas les recalculer à chaque image. (ex : approche des keyframes : calculer les features clés sur certaines images et les appliquer aux autres. Manque de précision. \footcite{zhou2022})
La segmentation sémantique reste toutefois lacunaire sur des vidéos longue durée, ou sur des vidéos dont les catégories sémantiques sont inconnues, ce qu'on apelle la segmentation "open world". \footcite{zhou2022} (les modèles actuels connaissent les catégories sémantiques à l'avance et ne peuvent pas gérer ou reconnaitre des catégorie inconnues). Zhou et al coseillent également d'adapter le calcul en fonciton du contenu de l'image, pour adapter des parties du réseau au contexte de l'image pour éviter un gaspillage de calcul. \footcite{zhou2022} 

Adapté au contexte skinsoft, il faudrait donc produire un mapping des champs utilisés pour décrire un film, une archive audiovisuelle, ou annoter des archives avec les champs dédiés. Les corpus sur lesquels sont entraînés les modèles sont bien loin de la complexité des archives audiovisuelles. 


%\subitem les méthodes orientées précision

%Ces méthodes agrègent l'information temporelle entre les images pour améliorer la prédiction. Elles utilisent le flux optique ou des modles basés sur l'attention. Ex de méthodes : TDNET \footcite{zhou2022}

%\subitem méthodes orientées vitesse

%Elles réutilisent les features des images adjacentes au lieu de les recalculer à chaque fois, ex : approche des keyframes : calculer les features clés sur certaines images et les appliquer aux autres. Manque de précision. \footcite{zhou2022}

Ex de résultats : Sur Cityscapes (VSS), la méthode EFC obtient les meilleures performances avec 83,5 pourcents d'IoUclass. \footcite{zhou2022}

%limites et manques : 
%Zhou et al identifient des secteurs manquants de la recherche sur la segmentation video sémantique : notamment la segmentation vidéo longue durée, la segmentation "open world" (les modèles actuels connaissent les catégories sémantqiues à l'avance et ne peuvent pas gérer ou reconnaitre des catégorie inconnues). Appliqué au contexte skinsoft, il faut faire un mapping des catégorie/champs.Les modèles atcuels sont entraînés sur des jeux de donénes qui comportent des catégories éloignées d'archives audiovisuelles. \footcite{zhou2022}
%Autre manque : solution d'anotations efficace : restent sous utilisées et explorées.\footcite{zhou2022}
%Il faudrait aussi adapter le calcul : adapter des parties du réseau selon le contenu de l'image pour éviter le gaspillage de calcul. \footcite{zhou2022}

Segmentation guidée par le langage (LVOS)

%revoir la différence entre segmentation sémantique, segmentation d'objets, segmentation temporelle : pas très clair%

%déjà intégrée dans la rédaction plus haut 
%Les diférents types d'algorithmes de segmentation qui peuvent être appliqués à la vidéo sont : l'identification des coupures de plans, la segmentation selon des critères visuels, la segmentation du son entre les paroles et la musique par exemple. \footcite{carrive2014}

problématiques liées à la nature de la vidéo : "C’est donc un problème de segmentation, qui est posé parce que les divisions "calculables"d’une vidéo (les plans et les segments de parole ou de musique), n’ont pas de valeur documentaire. En revanche, les divisions documentaires "naturelles" (sujets d’un journal, numéros d’unspectacle de variétés, scènes d’un film) ne sont pas directement calculables"\footcite{zotero-item-644}

%déjà rédigé plus haut
%La segmentation vidéo comprend à la fois la segmentation sémantique (objets dans l'image) et la segmentation temporelle, propre à la vidéo, qui segmente les différentes images dans le temps. 

%déjà rédigé plus haut
%La segmentation vidéo sémantique se divise en deux branches : segmentation d'objets vidéos, et segmentation sémantique.\footcite{zhou2022}

La segmentation d'objets vidéo se décline en trois mode : automatique, semi automatique, interactif. \footcite{zhou2022}

%déjà rédigé plus haut
%La segmentation d'objets vidéo (video foreground/background segmentation) consiste à segmenter les objets dominants dans l'image. \footcite{zhou2022}

(en quoi est ce différent d'une image fixe ? ça ne peut pas être appliqué aux images fixes ?). C'est plutot utilisé pour le montage : retirer des objets dans l'image par ex. Ca ne concerne pas la signification des images dans l'image, seulement leur place. 

%déjà mis dans la rédaction plus haut%
%La segmentation sémantique extraie des objets qui apparatiennent à de scatégories sémantqiues prédéfinies : "VSS aims to extract objects within predefined semantic categories (e.g., car, building, pedestrian, road) from videos" \footcite{zhou2022}
%Le but de la segmentation sémantique est d'assigner des étiquettes à chaque pixel extrait de l'image. 

%à rédiger 




\item pourquoi segmenter ?

La segmentation découle d'une volonté d'indexer le contenu de l'image. En effet, selon Carrive, "La première étape du processus d’analyse des documents audiovisuels est bien souvent une étape de segmentation temporelle qui consiste à identifier dans le document des séquences possédant des caractéristiques communes ou des ruptures entre ces séquences."\footcite{carrive2014}. Segmenter ainsi l'image permet d'indexer non plus la vidéo au complet mais des fragments, rendant son analyse plus précise. 
C'est également la manière dont sont analysés les films , plans par plans, ou séquences par séquence, pour en déceler le sens, à la fois une segmentation des éléments qui se trouve dans l'image, leur relations entre eux, et l'enchaîement des plans, les mouvements de caméra. L'entraînement de modèles pour l'indexation est donc calquée sur le regard humain qui permet d'analyser et d'indexr des films. 
Une fois la segmentation effectuée, l'objet d'analyse n'est plus le film entier mais des plans ou des séquences. \footcite{cordell} La segmentation crée ainsi de nouveaux objets indexables.
La segmentation est ainsi la premère étape de la compréhension d'un fichier vidéo : "One of the first tasks in understanding a video file is to break up the sequence of frames into distinct units" \footcite{arnold2019}
Une fois le fichier vidéo segmenté, il est alors possible d'organiser le contenu, d'améliorer l'accès, et de produire des analyses sur des larges données vidéo. \footcite{attard2025}
La segmentation permettrait d'analyser de larges quantités de données, tout en homogénéisant les pratiques documentaires, en augmentant la précision et la cohérence. \footcite{attard2025}.

%Proper segmentation facilitates content organisation, improves accessibility, and enhances the ability to perform analytics on large video datasets. »\footcite{attard2025}

%La segmentation comme logique documentaire : la segmentation d'une vidéo crée de nouveaux éléments indexables
%Cordell : "your object of study is no longer an entire film but shots or scenes" \footcite{cordell}

Qu'est-ce que ça signifie pour l'institution de passer d'un objet-film à une collection de plans ? Quelles sont les implications pour le droit, pour l'accès public, pour la description archivistique ?

%La segmentation découle d'une logique et d'un objectif d'indexation : problématique principale pour indexer : il faut déjà séparer la vidéo en tous les plans qui la composent pour en extraire les éléments : 
%"One of the first tasks in understanding a video file is to break up the sequence of frames into distinct units known as shots—or uninterrupted footage between cuts—in a process known as shot boundary detection."\footcite{arnold2019}

%Une vidéo doit être segmentée avant de pouvoir être indexée : 
%« One of the most common initial steps for video content organisation is segmenting the video into different meaningful parts. Proper segmentation facilitates content organisation, improves accessibility, and enhances the ability to perform analytics on large video datasets. »\footcite{attard2025}

%Pour justifier la mise en place d'une segmentation : « Automated segmentation accelerates the processing of large volumes of video content and ensures consistency and accuracy in the segmentation process. » \footcite{attard2025} 
la mise en place d'une automatisation de la segmentation permettrait d'homogénéiser l'indexation documentataire : un algorithme segmente les vidéos de la même façon, pas les humains. Mais les vidéos sont toutes différentes : est ce qu'on peut garder la même pratique de segmentation sur des corpus vidéo différents ?

\item Segmentation sémantique ou segmentation d'images
 
 La segmentation sémantique est une tâche de la vision par ordinateur. 
 
"Semantic segmentation is a fundamental task in computer vision. It aims to divide the image into multiple regions according to the class of the pixel, which is important for practical applications such as autonomous driving, automatic delivery sorting, and intelligent surveillance systems"\footcite{yuan2025}

Importance de la segmentation sémantique pour les vidéos : "videos contain not only spatial features but also temporal features. This naturally results in lower accuracy of image-based segmentation models compared with video-based segmentation models for Video Semantic Segmentation (VSS)." \footcite{yuan2025}

\item Pipelines de segmentation 

\subitem Pyscenedetect 
Une idée de pipeline de segmentation : est la détection de scènes avec pyscenedetect (découpage d'une vidéo en segments en détectant les coupes visuelle), puis classification des segments. D'un côté la segmentation syntaxique et de l'autre la classification sémantique. \footcite{attard2025}

\subitem ResNet (pas de la segmentation : créer un fichier sur la classification ?)
Le meilleur classifieur d'images avec le plus haut résultat de précision selon attard 2025 est ResNet
\footcite{attard2025}

Reccommandations : 
Les modèles multimodaux qui combinent audio et vidéo (vivit-ast) ne donnent pas les meilleurs résultats. Il faudrait donc adopter des architectures plus simples, qui demandent moins de calcul d'usage et qui séparent les pipelines. \footcite{attard2025}

\subitem ViViT
Expérimentation avec ViViT : 
le modele ViViT, modèle basé sur des transformers, capture les caractéristiques spatio temporelles de la vidéo, contrairement à Resnet qui analyse chaque image d'un vidéo indépendemment. ViViT extrait "extracts spatio-temporal tokens from the input video, which are then encoded by a series of transformer layers. » \footcite{arnab2021}. C'est un modèle qui tente d'analyser l'espace et le temps dans une vidéo donnée. Vivit adopte donc une approche "video-native". 
Le modèle à choisir dépend de la nature du corpus audiovisuel que l'on souhaite indexer. Si les archives sont des plans fixes qui ne changent que très peu, resnet suffit. 

Conclusion sur ViViT : 
« ViT has been shown to only be effective when trained on large-scale datasets, as transformers lack some of the inductive biases of convolutional networks, and thus require datasets larger than the common ImageNet ILSRVC dataset to train. » \footcite{arnab2021}. : les transformers n'ont pas les biais des CNN qui leur permettraient de faire des généralisations et de reconnaître un objet où qu'il soit dans l'image. Ils ont donc besoin d'être entrainés sur beaucoup plus de données : c'est une grande limite. Une approhe par calssifieurs d'images plus légers comme ResNet peut donc être préférable. ViViT obtient de très bons résultats mais entraînés sur des corpus vidéos contemporains : vidéos youtube. Diffiilement adaptable à un corpus d'archives historique sauf pré entraînement massif. 

\subitem 2SDS

l'étude de xin2023 \footcite{xin2023} présente 2SDS, un algorithme de segmentation temporelle vidéo. Plutôt que d’utiliser un réseau de neurones pour détecter les changements de scènes, 2SDS compare des empreintes (hash) d’images successives. Si deux frames sont suffisamment différentes (distance de Hamming supérieure à un seuil), elles appartiennent à des scènes distinctes.
« 2D CNNs only recognise a video as discrete images, rather than a continuous stream of images. This poses some issues. For example, a CNN model could not resolve the motion of a person because the person is stationary in every frame, and this will cause the loss of significant information in video analysis. »\footcite{xin2023} : un CNN 2D comme resnet analyse chaque plan de manière indépendante et ignore la logique temporelle de la vidéo. 2SDS résout ce problème en regroupant les plans similaires avant de les classifier. (à relier avec la notion de shot boundary detection). 
« We need to devise an implementation that complements the CNN model to solve the continuity issue. This implementation should group the discrete frames (adjacent on the temporal axis) that look similar to each other into scenes, this procedure is what we call temporal segmentation. »\footcite{xin2023} : la segmentation temporelle : regrouper les plans qui se suivent et visuellement similaires : c'est ce que fait pyscene detect avec le content detector. 
l'article explique aussi la mise en place d'un pipeline de détection de changement de plan : 
4 étapes : 
\subitem down sampling : miniaturisation du plan pour détecter les changements globaux, pas locaux
\subitem conversion de l'image RGB en niveaux de gris 
\subitem calcul d'une valeur de "hachage" de l'image : Chaque ligne de 9 pixels est convertie en 8 bits binaires comparant les pixels adjacents. Produit une « empreinte » numérique de la frame.
\subitem distance de hamming : 
comparaison des empreinte de deux frames consécutives. Si la distance est supérieure à 5, les deux plans apparatiennent à des scènes différentes. le seuil est paramétrable. pyscenedetect est l'implémentation concrète de ce processus. 

les auteurs conseillent de ne pas utiliser des réseaux de neurones pour la segmentation vidéo : « The usage of RNN or even neural networks is not practical in time sensitive tasks like real-time object recognition and live video stream analysis, which requires fast responding algorithms, and RNNs usually cannot meet those requirements. »\footcite{xin2023}

détecter les coupures entre es scènes est une analyse syntaxique et non sémantique : il suffit de détecter la différence visuelle. les meilleurs résutats de détection de coupures se font sur des plans fixes, avec des suets distants, des mouvements lents, docn plutot bien adaptées aux archives audiovisuelles. 

l'article de \cite{xin2023} ne cherche pas à comprendre les scènes mais seulement à analyser leur syntaxe pour en détecter les coupures. 2SDS règle le problème de délimitation en séquence, mais ne décrit pas la vidéo comme d'autres outils peuvent le faire : CLIP, YOLO, BLIP. De plus, le corpus de tets ets réduit (6 vidéos issues de youtube), et en couleur. 


\subitem video semantic segmentation (VSS) : TDNSNet

"fundamental machine vision task with various practical applications, such as autonomous driving and automated surveillance. Current studies mainly utilize temporal features based on optical flow and the self-attention mechanism to improve VSS accuracy."\footcite{yuan2025}
LImites : "reduced accuracy and computational overhead due to inaccurate optical flow and the computational cost of the self-attention mechanism."\footcite{yuan2025}
L'outil temporal difference segmentation net se propose comme uen solution à ces limites sellon l'article de yuan \footcite{yuan2025}

Limites identifiées par Yuan2025 des modèles de VSS existants : les modèles basés sur le flux optique (optical flow) sont dépendants de la prédiction du fluxet génèrent des erreurs, dégrandant ainsi la qualité de la segmentation (erro accumulation). Les modèles fondés sur les mécanismes d'auto attention capturent mieux les relations entre images mais leur complexité les rend difficilement applicables à des vidéos complexes elles mêmes.\footcite{yuan2025}

TDSNet : "TDSNet employs temporal-difference-based temporal feature through the Temporal Feature Refine Module (TFRM). To further improve the accuracy of VSS, TDSNet adaptively fuses temporal features of varied motion magnitude with the Motion Magnitude Refine Module (MMRM). This module weighs and fuses temporal features of different magnitudes between frames." \footcite{yuan2025}

TDSNet est donc un modèle qui se base sur la différence temporelle : 
"the temporal-difference-based features are obtained by the subtraction of the features from two consecutive frames. As a result, the model activates the dynamic regions of the feature while deactivating the static regions, allowing it to capture temporal features effectively." \footcite{yuan2025} Le modèle soustrait les features de deux images consécutives pour en faire ressortir les régions dynamiques (en mouvement) et désactive les régions statiques. Ainsi il capture uniquement l'information temporelle, pour un faible cout computationnel. 

TDSNet repose sur deux modules : 
TFRM : temporal feature refine model : il extrait les features temporelles par différence entre l'image cible (target frame) et l'image de référence (reference frame). Il utilise ensuite l'information de mouvement retenue pour aligner et fusionner les features de la référence dans l'image cible.\footcite{yuan2025} 

MMRM : Motion magnitude refine module 
les features extraites par le TFRM contiennent des magnitudes de mouvements variés : ""we observe that motion magnitudes in the far frames are not necessarily larger than that of near frames. For example, with repetitive motion of the video, farther frames may exhibit smaller motion magnitudes.". Le MMRM calcule donc le poids adaptatif de chaque feature temporelle selon la magnitude de mouvement réellement observée dans la vidéo en entrée, puis les fusionne.\footcite{yuan2025}

Le modèle est entraîné sur deux jeux de données : VSPW (13 images par secondes et un écart de mean Intersection over Union de 0.7 pourcent)


\subitem CamVid

Cambridge driving : base de données vidéo etiquetée par des classes sémantiques d'objets. \footcite{brostow2009} : 2009 donc trop vieux 
ce qui est intéressant, c'est que en 2009, il y a moins d'algorithmes puissants et précis qui existent, donc c'est la vé

\item Automatic temporal vidéo scene segmentation (ou video story segmentation)

"The automatic temporal video scene segmentation (also known as video story segmentation) is still an open problem without definite solutions in most cases." \footcite{trojahn2021}
les meilleurs solutions et qui donnent les meilleurs résultats qui existent sont les techniques mutlimodales utilisant les caractérstiques de plusieurs outils : "Among the available techniques, the ones which shows better results are multimodal using features extracted from multiple modalities."\footcite{trojahn2021}
La fusion multimodale peut se faire de deux façon, fusion en une seule représentation, fusion par segmentation de chaque modalité : "Multimodal fusion may be performed to fuse each modality as a single representation (early fusion) or by each modality segmentation (late fusion), the latter been widely due to multimodal fusion simplicity."\footcite{trojahn2021}. C'est la seconde approche qui est privilégiée. 

limites des CNN : "CNNs cannot adequately learn cues which are temporally distributed along the video due to difficulties to model temporal features data dependencies."\footcite{trojahn2021}

RNNs : "Successfully employed on text processing, RNNs are fitted to analyze sequences of data of variable length and may better grasp the temporal relationship among low-level features of video segments, hopefully obtaining more accurate scene boundary detection."\footcite{trojahn2021}. Les RNN permettent de mieux capturer les relations temporelles. 

Une des solutions serait de combiner CNN et RNN : complémentaires. 

\subitem DIVAN - INA
mise en place d'un système expérimental d'indexation semi automatique des archives audiovisuelles, conçu en trois étapes : segmenter, classifier, composer. \footcite{zotero-item-644}. les évènements détectés dans l'image sont annotés et validés par les archivistes.

\item segmentation audio 
la segmentation audio peut ête divisée en différentes sous tâches : 
la transcrpiption de la parole ("Les performances de ces systèmes peuvent être très bonnes lorsque de bonnes conditions sont réunies : bonnes conditions d’enregistrement, élocution claire, syntaxe correcte, noms propres connus du système en particulier." \footcite{carrive2014}), l'alignement audio/texte, indexation ("identification des termes les plus importants (mots clés), identification des sujets traités ou segmentation thématique" \footcite{carrive2014})

Une autre tâche de la segmentation audio et de l'analyse audio est l'identification des locuteurs. "Pour cela, il est nécessaire de posséder une référence pour chacun des locuteurs à reconnaître. Une référence peut consister en un ensemble d’échantillons de la voix de la personne à reconnaître"\footcite{carrive2014}

\item analyse des visages 
l'analyse des visages dans une vidéo se divise en plusieurs sous tâches : détection, suivi, caractérisation, reconnaissance. \footcite{carrive2014}

\item autres tâches de la segmantation 
Recherche d'instance (instance search) : "À partir d’un plan vidéo cible et d’un ensemble de thèmes qui délimitent la requête, il s’agit de retrouver dans une collection les plans les plus pertinents correspondant à la requête." \footcite{carrive2014}

Indexation sémantique : "Il s’agit de caractériser des plans vidéo en leur assignant des étiquettes (tags) représentant des caractéristiques abstraites (appelées « concepts »)." \footcite{carrive2014}

\item Découverte d'objets visuels 
"La découverte d’objets visuels consiste à rechercher automatiquement les objets apparaissant le plus fréquemment dans une collection d’images ou de vidéos" \footcite{carrive2014}

\item data mining 

"La fouille de données, ou data mining, recouvre un « ensemble de techniques ayant pour objet l’extraction d’un savoir à partir de grandes quantités de données, par des méthodes automatiques ou semi-automatiques » [14]. Les outils informatiques cherchent à découvrir, à faire émerger, à dégager, dévoiler des informations qu’il n’aurait pas été possible d’obtenir manuellement, parce que la quantité de documents à analyser est trop importante ou parce que les informations cherchées ne sont pas explicites et n’apparaissent en quelque sorte que statistiquement."\footcite{carrive2014}

Les archives audiovisuelles s'appliquent mal à ce concept puisque ce sont des archives : on veut pouvoir les décrire à la pièce, comme le ferait l'archiviste. One n eveut pas extraire de statistiques mais des descriptions précises. 

\item fossé sémantique 
désigne le fossé entre ce que les algorithmes savent détecter et ce que les utilisateurs recherchent. 



\end{itemize}

